\section{Introduction}
\label{sec:intro}

Modern large language models (LLMs) are typically deployed in an
autoregressive (AR) regime: tokens are emitted one at a time, and each new
token requires a forward pass through the full network. As models grow into
the tens of billions of parameters, this sequential dependence pushes
end-to-end decoding latency into the regime where wall-clock time is
dominated by memory bandwidth, not compute. Reducing this latency without
changing the output distribution is the central problem we study in this
work.

Inference is where the cost of LLMs is paid: coding copilots, chat
assistants, and RL rollouts all sit on the same critical path of
producing $T$ tokens from a target $M_t$ as quickly as possible while
preserving $p_{M_t}(\cdot)$. Any technique that reduces the number of
\emph{serial} forward passes through $M_t$ without biasing samples
directly improves throughput and user-perceived latency. Combining two
fast-moving lines of work, \emph{speculative decoding} and
\emph{discrete diffusion language modeling}, is therefore a path to
compounding speedups on real workloads.

\paragraph{Related work.}
\textbf{Speculative decoding} (SD)~\cite{leviathan2023, chen2023} couples a
fast AR draft model $M_d$ with the target $M_t$: $M_d$ proposes $k$ tokens
which $M_t$ verifies in a single parallel forward pass. Accepted tokens are
kept, and a corrected token is appended at the first rejection, preserving
$p_{M_t}$ exactly under standard rejection sampling. Speedups are
proportional to the acceptance rate but bounded by the fact that drafting
and verification run \emph{sequentially}. \textbf{Speculative$^2$ Decoding}
(S$^2$D)~\cite{saguaro2024} removes this idle time by having $M_d$
anticipate the most likely verification outcomes and pre-generate a tree of
$f$ candidate continuations \emph{in parallel} with the target's verify
step; if the verifier's actual decision matches one of the speculated
branches, the next block of drafts is already in hand. \textbf{Discrete
diffusion language models}~\cite{sahoo2024, arriola2025} take an orthogonal
view of generation: instead of producing tokens left-to-right, they
iteratively denoise a sequence of $[\textsc{mask}]$ tokens over a small
number of refinement steps, generating an entire block of tokens in
parallel. MDLM~\cite{sahoo2024} introduced a simple, scalable masked
diffusion objective, and BD3LM~\cite{arriola2025} extended this with a
block-causal (``staircase'') attention structure that gives the model a
left-to-right inductive bias while keeping inside-block parallelism.

\paragraph{Our approach.}
Existing SD and S$^2$D systems treat the draft as a black box AR
sampler, so they inherit the AR bottleneck inside the draft model itself.
We replace the AR drafter with a discrete diffusion drafter and run it
inside the S$^2$D loop. Concretely, at each draft step the diffusion model
fills $k$ masked tokens in $R$ parallel refinement steps ($R \ll k$), and
S$^2$D's lookahead tree is built on top of these diffusion-generated
blocks. Compared to vanilla SD, drafting itself is faster (parallel rather
than $k$-serial); compared to S$^2$D with an AR drafter, we shorten the
critical path of each speculation cycle. As a stretch extension, we also
experiment with a \emph{block-diffusion} drafter (BD3LM) that admits a
staircase attention pattern and prefix-KV reuse, which is closer in spirit
to a block-causal verifier.

\paragraph{Contributions.}
(i) We extend an open S$^2$D inference engine with a \texttt{block}
draft backend supporting MDLM-style and BD3LM-style discrete diffusion
drafters, three samplers (mask-predict, remask, first-hitting), and
two attention layouts (full, staircase) with optional prefix-KV reuse
(Sec.~\ref{sec:method}). (ii) We add a \emph{switching} drafter that
warm-starts with an AR drafter for the first $N$ completion tokens
before handing control to the diffusion drafter, so block diffusion is
only paid for once it has enough context. (iii) We benchmark against
AR, sync SD, and async S$^2$D on four datasets and compare against
vLLM and SGLang speculative decoding
(Sec.~\ref{sec:experiments}).
